% !TeX program = lualatex
% Compile twice: lualatex 223.tex
% All references and prose are embedded; no BibTeX database or external inputs.
\documentclass[11pt,a4paper]{article}
\usepackage[margin=24mm,headheight=14pt]{geometry}
\usepackage{fontspec}
\setmainfont{Latin Modern Roman}
\setsansfont{Latin Modern Sans}
\setmonofont{Latin Modern Mono}
\usepackage{amsmath,amssymb,mathtools}
\usepackage{unicode-math}
\setmathfont{Latin Modern Math}
\usepackage{microtype}
\usepackage{enumitem,xurl,needspace}
\usepackage[dvipsnames]{xcolor}
\usepackage{hyperref}
\hypersetup{unicode=true,colorlinks=true,linkcolor=MidnightBlue,urlcolor=MidnightBlue,
 pdftitle={Research Notebook: Section 223},pdfauthor={Research notebook},bookmarksnumbered=true}
\usepackage{bookmark}
\usepackage{tocloft}
\setlength{\cftsecnumwidth}{3em}
\cftsetpnumwidth{2.5em}
\cftsetrmarg{4em}
\usepackage{titlesec}
\titleformat{\section}{\Large\bfseries\raggedright}{\thesection}{0.7em}{}
\usepackage{fancyhdr}
\pagestyle{fancy}\fancyhf{}
\fancyhead[L]{\small\sffamily Open Problems | Research notebook}
\fancyhead[R]{\small 223 | 27 September 2026}
\fancyfoot[C]{\thepage}
\setlength{\parindent}{0pt}
\setlength{\parskip}{5pt plus 1pt}
\setlength{\emergencystretch}{3em}
\setcounter{tocdepth}{1}
\setcounter{secnumdepth}{1}
\setlist[itemize]{leftmargin=3em,itemsep=5pt,topsep=4pt,parsep=0pt}
\allowdisplaybreaks
\newcommand{\N}{\mathbb{N}}
\newcommand{\Z}{\mathbb{Z}}
\newcommand{\Q}{\mathbb{Q}}
\newcommand{\R}{\mathbb{R}}
\newcommand{\C}{\mathbb{C}}
\newcommand{\F}{\mathbb{F}}
\newcommand{\A}{\mathbb{A}}
\newcommand{\PP}{\mathbb P}
\newcommand{\E}{\mathbb{E}}
\newcommand{\eps}{\varepsilon}
\newcommand{\dd}{\,\mathrm d}
\newcommand{\sref}[2]{\hyperlink{source:#1:#2}{[S#2]}}
\newcommand{\eref}[2]{\hyperlink{extra:#1:#2}{[E#2]}}
% Turn the next switch off for a shorter reading copy. The quotations preserve
% every exactTarget field, including qualifications omitted from a short summary.
\newif\ifshowcatalogue
\showcataloguetrue
\newcommand{\cataloguescope}[1]{\ifshowcatalogue
 \paragraph{Exact scope in the supplied catalogue.}
 \begin{quote}\small #1\end{quote}\fi}

% Independently compilable extraction; original section text retained.
\numberwithin{equation}{section}
\begin{document}
\setcounter{section}{222}
% ================================================================
% problemId: problem.babais-diameter-conjecture-for-finite-simple-groups
% source releaseRank: 223; source releaseStatus: open_with_solved_subcases
% exactTarget SHA-256: 19b8391fa2e93e86ea1ffb58f43afaf33004dd9d1ddcb77f25b11879c97258c8
\clearpage
\section{\texorpdfstring{Babai's diameter conjecture for nonabelian finite simple groups}{Babai's diameter conjecture for nonabelian finite simple groups}}\label{223}
\subsection{Definitions and mathematical statement}
A nonabelian finite simple group has no nontrivial proper normal subgroup and is not commutative. For a generating set $X$, its undirected Cayley graph joins $g$ to $gx$ for $x\in X\cup X^{-1}$. Its diameter is the largest, over $g\in G$, of the shortest word length representing $g$ in these steps. The conjecture asks for one absolute $c>0$ such that
\[
 \operatorname{diam}\operatorname{Cay}(G,X\cup X^{-1})\le(\ln|G|)^c
\]
for every such $G$ and every generating set. Constants and the exponent may not depend on the group family, rank, or generators. Results requiring a special element in $X$ are subcases rather than the full assertion.
\par\smallskip\textit{Formulation and concept references:} \sref{223}{1}.
\cataloguescope{Is there an absolute constant c>0 such that for every nonabelian finite simple group G and every generating set X of G, the undirected Cayley graph with steps X union X\textasciicircum{}\{-1\} has diameter at most (ln |G|)\textasciicircum{}c? Thus every element of G is a word of that length in X and its inverses, with the same c for all groups and generating sets.}
\subsection{Short English statement}
Can every element of every nonabelian finite simple group be expressed in any generating set using only a polylogarithmic number of generators and inverses?
% BEGIN RESEARCH ATTEMPT 223
\subsection{Research attempt}

\paragraph{Current frontier.}
Bounded-rank product theorems yield polylogarithmic diameter with an
exponent depending on rank. The source also proves a uniform exponent
for specified classical groups when a transvection is already in the
generating set, retaining its odd-field and exceptional-field conditions.
\sref{223}{1}, Theorem~1.2, Corollary~1.3 and Theorem~1.5.
The 2024 work completing the Liebeck--Nikolov--Shalev conjecture gives a
short product of \emph{conjugates} of an arbitrary subset, but does not
bound the original-generator word lengths of the conjugating elements.
\eref{223}{1}, Abstract. That difference is central to an attempted
transfer to this diameter question. Neither result supplies a uniform
answer for every generating set of every simple group.

\paragraph{Attack route.}
A rank-independent product-growth exponent would immediately control
iterated balls, but is false in that naive form. Product decomposition
by conjugates is more flexible, yet may conceal long conjugators.
A third route first manufactures a useful seed element, such as a
transvection, by a short word, and then invokes a known special-element
diameter theorem. We test all three issues explicitly. The selected
construction below disproves the naive growth lemma, not Babai's bound.

\paragraph{Attempt.}
\textbf{1. The exact cost of growth iteration.}
Put $A=X\cup X^{-1}\cup\{1\}$. Suppose for a moment that every
intermediate generating set $B$ satisfies either $B^3=G$ or
$|B^3|\ge|B|^{1+\eta}$, with one $\eta>0$. After $j$ iterations,
$|A^{3^j}|\ge|A|^{(1+\eta)^j}$ until coverage occurs. Taking the first
$j$ for which the latter exceeds $|G|$ shows
\begin{equation}\label{223:iteration}
 \operatorname{diam}(G,X)\le
 3\left(\frac{\log|G|}{\log|A|}\right)^{\log3/\log(1+\eta)}.
\end{equation}
Thus an exponent $\eta$ tending to zero with rank loses the desired
uniformity. This reproduces the quantitative reason for the rank
restriction in \sref{223}{1}, Corollary~1.3; it does not improve that
corollary.

\textbf{2. An explicit family rules out uniform tripling growth.}
For $k\ge4$, take $G=A_{2k+1}$, let $H=A_k$ act on the first $k$
letters, and let $g=(1\ 2\ \cdots\ 2k+1)$. This odd-length cycle is
an even permutation. Set
\[
 A=H\cup\{g,g^{-1}\}.
\]
It is symmetric, contains the identity and generates $G$. Indeed $H$
contains $(1\ 2\ 3)$, whose conjugates by powers of $g$ give consecutive
three-cycles; these generate the alternating group. One can see the last
claim inductively: $A_{m-1}$ together with $(m-2\ m-1\ m)$ generates
all three-cycles involving $m$, by conjugation, and hence $A_m$; the small
initial cases are checked directly.

The intersection $H\cap gHg^{-1}$ is the alternating group on the
$k-1$ common letters, so
\[
 |HgH|=k|H|,\qquad |Hg^{-1}H|=k|H|.
\]
Words of length three in $A$ contain zero, one, two, or three letters
from $\{g,g^{-1}\}$. The one-letter cases lie in these two double
cosets. The two-letter cases are among four sets $g^{\pm1}Hg^{\pm1}$
and the at most six left or right cosets obtained from $g^0,g^2,g^{-2}$.
The zero-letter case is $H$, and the three-letter case has at most four
values. The deliberately generous bound
\begin{equation}\label{223:smalltripling}
 |A^3|\le(2k+12)|H|+4
\end{equation}
follows. It is much smaller than $|A_{2k+1}|$ for large $k$.
Since $|A|=k!/2+2$,
\[
 \frac{\log|A^3|}{\log|A|}
 \le 1+O\left(\frac{\log k}{\log(k!)}\right)=1+O(1/k).
\]
Therefore there is no absolute positive $\eta$ making the proposed
tripling alternative valid for all finite simple groups and all
symmetric generating sets. This refutation is entirely compatible with
polylogarithmic diameter: a ball may grow slowly at one scale and later
grow much faster. It supplies no large-diameter counterexample.

\textbf{3. Short conjugates exist; short products in them remain an issue.}
Let $G$ be any nonabelian finite simple group and $a\ne1$. Write
$H_j$ for the subgroup generated by $bab^{-1}$ with $b$ a word of
length at most $j$ in $A$. If $H_j=H_{j+1}$, conjugating its generators
by $A$ shows that $H_j$ is normal in $G$. It contains $a$, so simplicity
gives $H_j=G$. Every strict inclusion at least doubles subgroup order.
Thus $H_j=G$ for some $j\le\lceil\log_2|G|\rceil$.
If $a$ itself has $X$-word length $L$, these conjugates all have length
at most $2\lceil\log_2|G|\rceil+L$.
This is the elementary normal-closure principle underlying
\sref{223}{1}, Lemma~2.1.

However, saying that these short conjugates \emph{generate} $G$ does
not bound the number needed in an expression for each element.
Conversely, the product-of-conjugates theorem in \eref{223}{1} bounds
the number of factors but allows conjugators anywhere in $G$.
Replacing those conjugators by words of length bounded by the unknown
diameter gives an inequality of the form
$D\le C(\log|G|)(2D+1)$, not an upper bound on $D$.
Combining the two statements in this circular manner does not close the
gap. A useful new lemma would have to keep both the number of factors
and the conjugator word lengths small simultaneously.

\textbf{4. A seed-element route and its missing quantitative input.}
In a special linear group, write rank-one unipotents as
$t=I+u\varphi$ and $s=I+v\psi$, where
$\varphi(u)=\psi(v)=0$. If also $\psi(u)=0$ and
$a=\varphi(v)\ne0$, direct multiplication gives
\begin{equation}\label{223:commutator}
 [t,s]=tst^{-1}s^{-1}=I+a u\psi.
\end{equation}
Indeed $U=u\varphi$ and $V=v\psi$ have $U^2=V^2=VU=0$,
$UV=a u\psi$, and all higher terms appearing in the multiplication
vanish. If $u,\psi\ne0$, the result is again a transvection.
The word length of the commutator is at most twice the sum of the two
input word lengths. This is a usable propagation formula, but it starts
with suitable unipotents and incidence conditions; it does not create
the first such element from arbitrary generators.

For any classical-group case covered by \sref{223}{1}, Theorem~1.5,
a transvection of word length at most $(\log|G|)^b$, for a universal
$b$, would finish that case: append the element to $X$, use the theorem,
and replace each appended letter by its original word. The exponents
add, and universal multiplicative constants can be absorbed into a
larger exponent because $|G|\ge60$. This retains the source's field
exceptions and does not address other families automatically.
No such short-seed theorem for every generating set is established here.
Random-generation results cannot be promoted to that universal input.

\paragraph{Stress test.}
The generating set in \eqref{223:smalltripling} is allowed to be large;
the original conjecture quantifies over all generating sets. The
construction disproves only a \emph{sufficient strengthening}, not the
target. The bounded-rank theorem does not apply with a uniform constant
to alternating groups of increasing degree. The normal-closure chain
has a logarithmic number of strict subgroup inclusions, not a logarithmic
Cayley diameter. A compressed straight-line description built from
commutators may expand to a long ordinary word; the target charges the
expanded word length.

The source's unitary field notation, odd-characteristic conditions and
exceptional field sizes are not silently discarded in the seed reduction.
The unrestricted finite-simple-group assertion therefore still needs
more than the transvection calculation.

\paragraph{Outcome.}
\textbf{Outcome: Exploratory approach with unresolved gap.}

The attempt constructs a generating family that invalidates the most
naive rank-independent tripling route, and gives explicit word-length
accounting for conjugation and commutator approaches. It isolates a
short-seed sufficient condition in the source's classical-group cases.
No priority claim is made for the normal-closure lemma or elementary
rank-one matrix identities.

The unresolved step is simultaneous control of growth, useful seed
elements and conjugating word lengths across all groups and generating
sets. No universal diameter estimate, and no family violating one, has
been proved.
% END RESEARCH ATTEMPT 223
\subsection{Sources}
\begin{itemize}\raggedright
\item[\textnormal{[S1]}]\hypertarget{source:223:1}{} Martino Garonzi, Zoltán Halasi and Gábor Somlai, On the diameter of Cayley graphs of classical groups with generating sets containing a transvection; arXiv2203.03323v1 March7,2022.
\url{https://arxiv.org/pdf/2203.03323v1}.
{\small\textit{Catalogue formulation source.}}
% BEGIN ADDED REFERENCES 223
\item[\textnormal{[E1]}]\hypertarget{extra:223:1}{}
N. Lifshitz, \emph{Completing the proof of the Liebeck--Nikolov--Shalev
conjecture}, arXiv:2408.10127v2, 26 September 2024, Abstract and the
product-of-conjugates theorem in the Introduction.
\url{https://arxiv.org/abs/2408.10127v2}. Consulted 27 September 2026.
The stated short product allows unrestricted conjugators; a word-length
bound for them is not inferred here.
% END ADDED REFERENCES 223
\end{itemize}

\end{document}
